\documentclass[12pt,a4paper]{article}
% \usepackage[utf8]{inputenc} % fontspecと併用時は不要
\usepackage[japanese]{babel}
\usepackage{amsmath}
\usepackage{amsfonts}
\usepackage{amssymb}
\usepackage{geometry}
\usepackage{graphicx}
\usepackage{enumitem}
\usepackage{fancyhdr}
\usepackage{verbatim}

% 日本語改行設定
\usepackage{xeCJK}
\usepackage{indentfirst}
\setlength{\parindent}{1em}

% 日本語フォント設定
\usepackage{fontspec}
\setmainfont{Hiragino Mincho Pro}
\setsansfont{Hiragino Sans}
\setmonofont{Menlo}
\setCJKmainfont{Hiragino Mincho Pro}
\setCJKsansfont{Hiragino Sans}
\setCJKmonofont{Menlo}
% ページ設定
\geometry{margin=1.5cm}
\pagestyle{fancy}
\fancyhf{}
\fancyhead[L]{プログラミング応用 第一回}
\fancyhead[R]{演習問題}
\fancyfoot[C]{\thepage}
\renewcommand{\headrulewidth}{0.4pt}
\renewcommand{\footrulewidth}{0.4pt}

% 問題番号のカウンター
\newcounter{problem}
\setcounter{problem}{0}

% シンプルな問題環境
\newenvironment{problem}[1][]{
    \refstepcounter{problem}
    \vspace{1em}
    \noindent
    \textbf{問\arabic{problem}} #1
    %\vspace{0.5em}
    %\par
}{
    \vspace{1em}
}

\begin{document}

% ヘッダー情報
\begin{center}
    {\Large \textbf{プログラミング応用 第一回}}\\[0.5em]
    {\large \textbf{演習問題}}\\[1em]
\end{center}

\noindent\textbf{注意事項}

\begin{itemize}[leftmargin=1em]
    \item 解答はpdf形式でLMSにアップロードしてください (例えば紙に手書きしてスキャンしたものや \TeX をコンパイルしたものなど)
    \item 締切：2025年10月14日（火）日本時間13時
    \item 質問がある場合は手を挙げてください。
\end{itemize}

\begin{problem}
オーダー記法の厳密な定義に基づいたとき、以下の主張はそれぞれ正しいだろうか？正しい場合は証明し、正しくない場合は反例を示せ。
\begin{enumerate}
    \item $3n^2+5n+2=O(n^2)$
    \item $f(n)=O(g(n))$かつ$g(n)=O(h(n))$ならば$f(n)=O(h(n))$
    \item $\log_2 f(n) = O(\log_2 g(n))$ならば$f(n)=O(g(n))$
\end{enumerate}
\end{problem}

\begin{problem}
以下のコードは何を計算しているか答えよ。また、その時間計算量も答えよ。ただし、計算量は演算の回数だけ考えればよく、答えだけオーダー表記で記述せよ。なお、コード内のlen関数の計算量は無視してよい。入力AとBはどちらも $n\times n$ の二次元リストである。

\begin{verbatim}
def f(A, B):
    n = len(A)
    C = [[0]*n for _ in range(n)]
    for i in range(n):
        for j in range(n):
            for k in range(n):
                C[i][j] += A[i][k] * B[k][j]
    return C
\end{verbatim}
\end{problem}

\begin{problem}
フィボナッチ数列$(a_n)_{n\ge 0}$に関する以下の事実を証明せよ。ただし、$a_0=0, a_1=1$とする。
\begin{align*}
  \begin{pmatrix}
    a_{n+1} \\
    a_{n}
  \end{pmatrix}
  =
  \begin{pmatrix}
    1 & 1 \\
    1 & 0
  \end{pmatrix}^n
  \begin{pmatrix}
    1 \\
    0
  \end{pmatrix}
\end{align*}

\end{problem}

\begin{problem}
自然数 $ n$ に対し, $n$ を一つ以上の自然数の和で表す方法(これを$n$の分割という)の総数を$p(n)$とする。例えば
\begin{align*}
n &= 4 \\
&= 3+1 \\
&= 2+2 \\
&= 1+1+2 \\
&= 1+1+1+1
\end{align*}
より$p(4)=5$である。また、$n$の分割のうち、$k$以上の自然数のみを使って表す方法の総数を$S(n,k)$とする。例えば
\begin{align*}
4 &= 3+1 \\
&= 2+2
\end{align*}
より$S(4,2)=2$である。$k<n$のときの$S(n,k)$を$S(n-k,k)$と$S(n,k+1)$を用いて表せ。
\end{problem}

\end{document}
